% ======================================================================
% SECTION 1: INTRODUCTION
% Three thematic blocks (a) FDA RTCT, (b) AI baseline, (c) Transition.
% ======================================================================

\subsection{The FDA Real-Time Clinical Trials Announcement}
\label{sec:intro-fda}

On 28 April 2026 the U.S. Food and Drug Administration announced two
major steps to advance the implementation of real-time clinical trials.
The agency unveiled the successful initiation of two proof-of-concept
trials that report endpoints and safety signals to the FDA in real
time, and released a Request for Information regarding a proposed pilot
program that will launch in summer 2026 \cite{fda2026realtime}. The
press release framed early-phase clinical trials as a bottleneck in
drug development, characterized by high uncertainty, limited patient
populations, and inefficient sequential decision-making. Data is
typically reported from sites to sponsors, who analyze and subsequently
submit data to the FDA, with the lag between site, sponsor, and agency
sometimes measured in months or years.

FDA Commissioner Marty Makary, M.D., M.P.H. described the change in
direct terms: \emph{"For 60 years, we've been conducting clinical trials
in the same way, where key data signals can take years to reach the
FDA. The lag time can delay regulatory decisions unnecessarily and slow
down the drug development timeline. We are boldly advancing a modern
approach whereby FDA scientists can view safety signals and endpoints
in real time as a trial progresses. This will help us accelerate
promising therapies, and build toward our ultimate goal of running
real-time, continuous trials across all phases of drug development."}
The Commissioner's framing positions the agency as enabling a faster
path from a promising therapy to a patient who needs it, not as the
limiting factor in that path.

The two proof-of-concept trials selected for the launch are
representative of distinct oncology contexts. AstraZeneca's TRAVERSE
is a Phase 2 multi-site trial in patients with treatment-naive mantle
cell lymphoma, with participation from The University of Texas MD
Anderson Cancer Center and the University of Pennsylvania. Amgen's
STREAM-SCLC is a Phase 1b trial in patients with limited-stage small
cell lung carcinoma, with final site selection in process at the time
of the announcement. The FDA met with each sponsor to establish
criteria for reporting signals in real time, and validated the
technical framework for AstraZeneca's trial through Paradigm Health.
Paradigm Health acts as the validated technical conduit between the
sites, the sponsor, and the FDA real-time stream.

Chief AI Officer Jeremy Walsh stated: \emph{"Real-time trials have been
talked about for years. We demonstrated that it is not only possible,
but also potentially transformative for the clinical trials ecosystem.
We have to consider our processes from the standpoint of a patient
awaiting a potentially powerful treatment."} The Request for
Information seeks community input on pilot program design, evaluation
metrics, and success criteria, with the comment period closing on 29
May 2026, final selection criteria disseminated in July 2026, and
pilot selections completed in August 2026.

The FDA framed real-time trials as an important step toward facilitating
\emph{continuous} trials. At present, most clinical development occurs
in discrete phases, run according to a protocol and typically as a
separate study. There is generally a hiatus in the development program
after one phase ends and the next begins, slowing the pace of product
development. Because real-time trials allow the FDA to view key
insights as the trial progresses, this hiatus could be eliminated or
reduced to a minimum, enabling continuous trials. This continuous
framework is the regulatory stage on which the four simulations in
this paper operate.

\subsection{Current AI Patient Prediction Capabilities}
\label{sec:intro-ai-baseline}

The current oncology AI patient-prediction landscape is heterogeneous,
narrow, and largely retrospective. The architectural reality is that
today's deployed prediction systems are not general-purpose reasoning
agents, but rather a collection of task-specific supervised models
fitted to fixed feature sets. Common architectures include ridge Cox
regression, support vector machines, random survival forests, gradient
boosted trees such as XGBoost, and elastic net classifiers
\cite{Smiley2025OncologyMLReportingQuality}. Methodological and
reporting quality remain mixed; many published models report a single
internal validation cohort with no external test, no fairness audit,
and limited prospective signal. The Huang 2025 systematic review and
meta-analysis of machine learning versus traditional Cox regression for
survival prediction in cancer using real-world data found a
standardized mean difference of 0.01 with a 95\% confidence interval
of -0.01 to 0.03, a null result that should temper any claim that
machine learning per se delivers a routine survival-prediction
advantage on structured tabular features
\cite{Huang2025CancerSurvivalMetaAnalysis}.

Short-horizon mortality prediction from electronic health record data
is the most mature application of supervised oncology AI. Manz 2020
validated a gradient boosting model to predict 180-day mortality for
24,582 outpatients with cancer, reporting an area under the receiver
operating characteristic curve of 0.89, a positive predictive value of
45.2 percent at the 40 percent decision threshold, and a 4 to 8 day
prospective lead time over routine clinical assessment
\cite{Manz2020Oncology180DayMortality}. Sidey-Gibbons 2022 showed that
patient-reported-outcome-only models can reach AUROC values of 0.69 to
0.76 for 180-day mortality in women with ovarian cancer, evidence that
even sparse non-clinical features carry useful prognostic signal
\cite{SideyGibbons2022OvarianPROMortality}. The ceiling for these
short-horizon EHR-based systems sits in the 0.80 to 0.89 AUROC range
across cohorts, and the ceiling has been stable for several years.

Adverse-event and hospitalization prediction during cancer treatment
is the application area in which supervised AI has produced its
strongest prospective evidence. Hong 2020 published the SHIELD-RT
prospective randomized study, which used a machine learning model to
direct twice-weekly clinical evaluations during radiation and
chemoradiation, and reported a 45 percent relative reduction in
acute-care events and a 48 percent reduction in associated costs
\cite{Hong2020SHIELDRT}. Elia 2025 reported a multi-institutional
external validation of SHIELD-RT across approximately 22,000 courses,
with AUROC values of 0.756 to 0.770 across the validation institutions
\cite{Elia2025SHIELDRTExternalValidation}. Li 2022 used XGBoost to
predict short-term cardiotoxicity in colorectal cancer patients
starting fluoropyrimidine-based chemotherapy at AUC 0.816
\cite{Li2022FluoropyrimidineCardiotoxicity}. Cai 2024 reported that
treatment-relevant toxicity in breast cancer patients undergoing
neoadjuvant systemic treatment can be predicted at AUROC 0.75 when
dose-intensity features are added \cite{Cai2024BreastCancerToxicityPrediction}.
Ugwu 2026 reported a 2026 systematic review and meta-analysis of head
and neck radiation toxicity AI models with a pooled AUROC of 0.76
\cite{Ugwu2026HeadNeckRadiationToxicityMeta}.

Medium and long-horizon survival prediction has matured through
multimodal weakly supervised models. Li 2025 reported AIM-LCpro for
non-small cell lung cancer with a C-index of 0.785 to 0.804 on internal
test data and 0.693 to 0.749 on external cohorts
\cite{Li2025AIMLCpro}. Yuan 2025 reported PROGPATH, a pancancer Vision
Transformer foundation model for prognosis, with C-indices of 0.713 to
0.805 across 17 external cohorts, the broadest pancancer evaluation
reported to date \cite{Yuan2025PROGPATH}. The 2025 image-based lung
cancer meta-analysis reported a pooled sensitivity of 0.83, a pooled
specificity of 0.83, and a pooled AUC of 0.90 across image-only
classifiers \cite{Yuan2025LungImageMeta}.

Response and progression prediction under immune checkpoint inhibition
is the application area with the most active multimodal research.
Yoo 2025 published SCORPIO, a routine-blood-test-and-clinical-data model
with a median time-dependent AUC of 0.763 on hold-out and 0.725 on
external validation \cite{Yoo2025SCORPIO}. Vanguri 2022 integrated
radiology, pathology, and genomics for response to PD-(L)1 blockade in
non-small cell lung cancer at AUC 0.80
\cite{Vanguri2022MultimodalPDL1NSCLC}. Captier 2025 reported a
late-fusion model that integrates clinical, pathological, radiological,
and transcriptomic data with a C-index of 0.75 for overall survival
and AUC 0.81 for one-year death \cite{Captier2025MultimodalNSCLCImmunotherapy}.
Arango-Argoty 2025 reported PBMF, a contrastive learning approach to
predictive biomarker discovery designed to improve clinical trial
outcomes \cite{ArangoArgoty2025PBMF}. Xu 2025 reported a multimodal
fusion system in unresectable hepatocellular carcinoma
\cite{Xu2025MMF}. Zhao 2023 reported anti-PD-1 radiomics in advanced
breast cancer at AUC 0.86 on internal test \cite{Zhao2023BreastRadiomics}.
Iivanainen 2021 used XGBoost on electronic patient-reported outcomes
to predict immune-related adverse events
\cite{Iivanainen2021irAE}. Brouwer 2025 used smartphone step count
data to predict clinical outcomes in patients with cancer receiving
systemic treatment \cite{Brouwer2025StepCount}. Sun 2024 modeled
patient health signals to inform long-term survival on immune
checkpoint inhibitor therapy \cite{Sun2024HealthSignals}.
Seyednasrollah 2017 organized a DREAM challenge to build prediction
models for short-term discontinuation of docetaxel in metastatic
castration-resistant prostate cancer, and Deng 2020 extended the
challenge with a treatment stratification model in the same disease
\cite{Seyednasrollah2017Docetaxel, Deng2020mCRPC}.

The reporting and methodological floor is set by two recent guidance
papers. Collins 2024 published the TRIPOD+AI statement, an updated
guidance for reporting clinical prediction models that use regression
or machine learning methods \cite{Collins2024TRIPODAI}. El Emam 2024
published the Consolidated Reporting Guidelines for Prognostic and
Diagnostic Machine Learning Models (CREMLS), which extend the TRIPOD+AI
recommendations toward operational deployment \cite{ElEmam2024CREMLS}.
Both standards apply to any new predictive system, including the
LLM-based simulations in this paper, that is intended for clinical
deployment.

The operational gaps left by this baseline are the subject of the rest
of this paper. The supervised models named above are retrospectively
trained, narrow per task, restricted to homogeneous endpoints inside
each task, and not coupled to a real-time signal-sharing channel
between sites, sponsors, and the FDA. They cannot ingest a heterogeneous
mix of protocol text, EHR records, robotic telemetry, regulatory
adaptations, and ASCII facility diagrams in a single inference pass.
They cannot emit a structured commit at hourly cadence to a public
GitHub repository. The four author simulations in this paper address
each of these gaps using a different architectural premise: a single
Claude Code Opus 4.7 Max instance with a 1M token context window holds
the entire trial repository in context and emits hour-by-hour patient
predictions, safety signals, and forecasts.

\subsection{Transition to the Four Simulations}
\label{sec:intro-transition}

Repository-scale processing with a 1M token context window is the
computational signature that separates this work from the supervised
baselines named in \S\ref{sec:intro-ai-baseline}. Claude Code Opus 4.7
Max can hold the protocol, regulatory, robotic, EHR, and ASCII-diagram
content of an entire oncology trial in a single inference pass and
emit hour-by-hour patient predictions, safety signals, and forecasts
at a cadence and depth that the narrow supervised models in current
oncology trial practice have not demonstrated
\cite{claude-opus-47, claude-code}. The thesis is computational and
contextual, not a head-to-head AUROC benchmark on a fixed test cohort;
the comparisons in \S\ref{sec:discussion} are scoped accordingly.

This paper documents four author Physical AI oncology trial
simulations. Simulation 1 is the Continuous National 24/7 RTCT, a
clinical trial site simulation across four sites (Houston,
Philadelphia, Boston, Texas Medical Center) and 116 robot instances,
with hour-by-hour text artifacts in
\path{new-trial/national-24-7-trial/hour-00} through \path{hour-55} and
the extra-hours block \path{extra-hours/hour-56} through \path{hour-83}
that contains approximated diagrams \cite{kawchak_2026_19176370,
kawchak_2026_19244918}. Simulation 2 is the Single-Patient 10-Stage
Journey, a clinical trial site simulation that follows patient
PAT-2026-0042 through ten Python-orchestrated stages from
pre-screening to closeout
\cite{kawchak_2026_19119939, kawchak_2026_18810541}. Simulation 3 is
the 24-Hour Autonomous Sponsor, a clinical trial sponsor simulation
that runs 53 core agents organized into four governance, study
execution, site/robotics, and trust layers \cite{kawchak_2026_19396256}.
Simulation 4 is the 168-Hour 7-Day Sponsor Extension, a clinical trial
sponsor simulation that produces 525 ASCII diagrams, 168 Python
scripts, 168 JSON outputs, and a complete local-verification run on a
2015 Intel Core i5-6200U laptop with 4 GB of RAM
\cite{kawchak_2026_19396256}. The on-screen Table of Contents follows
directly after this introduction on a fresh page.
